% Regression: hard10_v3 — imported current-behavior fixture from the handoff corpus.
% Formula: (gauge-fixed, A_L left-canonical, A_R right-canonical):
% hard10 v3 — tangent-space projector of a uniform MPS (RMP fig2-pepe_PTA),
% re-attempted with the upgraded keys: `nopair` rows, `west=cup` /
% `east=cup` cups (the tangent-space brackets), and `tri=l` / `tri=r`
% canonical triangle skins.
%
% Formula (gauge-fixed, A_L left-canonical, A_R right-canonical):
%   P_{T(A)} = sum_{i=1}^n  P_L^{(i-1)} (x) 1_i (x) P_R^{(i+1)}
%            - sum_{i=1}^n  P_L^{(i)}          (x) P_R^{(i+1)},
% with block projectors (orthonormal by the canonical conditions)
%   P_L^{(k)} = sum_alpha |Phi^L_alpha><Phi^L_alpha|,   A_L-chain, sites <= k;
%   P_R^{(k)} = sum_beta  |Phi^R_beta ><Phi^R_beta |,   A_R-chain, sites >= k.
%
% Key mapping per block: an all-legs-open ket-over-bra window
% (rows={ket:nopair, bra}) of canonical triangles; the bracket tying ket to
% bra at the block edge is `east=cup` (left blocks) / `west=cup` (right
% blocks); \tndots marks the semi-infinite continuation.  Pictures set side
% by side are tensor factors, so each summand is assembled exactly as the
% formula factors it: P_L (x) 1_i (x) P_R.
%
% Noted gap: the identity on site i (the RMP's bare blue vertical wire
% between the brackets) has no atom — every tenkz picture ink is anchored on
% a glyph, and a cup cannot sit at an interior hole edge (\tnskip leaves the
% gap's indices OPEN by design; the bracketed summand needs them TIED).  The
% factor 1_i is therefore carried by algebra, which contributes no drawn
% legs — hence both summands' DRAWN signatures match (see below).
%
% Signature balance across the minus sign (drawn ink, dots elide the rest):
%   sum 1: left block (vw 2, ve 0, pu 2, pd 2) + right block (0,2,2,2)
%          + algebraic 1_i (no ink)            = (2,2,4,4)
%   sum 2: left block (2,0,2,2) + right block (0,2,2,2) = (2,2,4,4)  — equal.
\documentclass[varwidth=19cm, border=6pt]{standalone}
\usepackage{amsmath, amssymb}
\usepackage{tenkz}
\begin{document}
\begin{align*}
  P_{T(A)} \;=\;& \sum_{i=1}^{n}\;
  % P_L^{(i-1)}: A_L window, all legs open, bracket (cup) at the east edge
  \begin{tenkz}[rows={ket:nopair, bra}, east=cup]
    \tndots & \tn[tri=l]{} & \tn[tri=l]{}\\
    \tndots & \tn*[tri=l]{} & \tn*[tri=l]{}
  \end{tenkz}
  \otimes \mathbb{I}_i \otimes    % identity at site i (no bare-wire atom)
  % P_R^{(i+1)}: A_R window, bracket (cup) at the west edge
  \begin{tenkz}[rows={ket:nopair, bra}, west=cup]
    \tn[tri=r]{} & \tn[tri=r]{} & \tndots\\
    \tn*[tri=r]{} & \tn*[tri=r]{} & \tndots
  \end{tenkz}
  \\[0.4em]
  -\;& \sum_{i=1}^{n}\;
  % P_L^{(i)} (x) P_R^{(i+1)}: the two brackets meet back to back
  \begin{tenkz}[rows={ket:nopair, bra}, east=cup]
    \tndots & \tn[tri=l]{} & \tn[tri=l]{}\\
    \tndots & \tn*[tri=l]{} & \tn*[tri=l]{}
  \end{tenkz}
  \otimes
  \begin{tenkz}[rows={ket:nopair, bra}, west=cup]
    \tn[tri=r]{} & \tn[tri=r]{} & \tndots\\
    \tn*[tri=r]{} & \tn*[tri=r]{} & \tndots
  \end{tenkz}
\end{align*}
\end{document}
